%==============================================================================
% ElegantBook PDF Style - LaTeX Preamble
% Replicates the visual style of ElegantBook template for Quarto PDF output
%==============================================================================

%------------------------------------------------------------------------------
% 1. Package Loading
%------------------------------------------------------------------------------
\usepackage{geometry}
\usepackage[table]{xcolor}
\usepackage[center]{titlesec}
\usepackage{fancyhdr}
\usepackage[most]{tcolorbox}
\usepackage[shortlabels,inline]{enumitem}
\usepackage{tikz}
\usetikzlibrary{backgrounds,calc,shadows,positioning,fit}
\usepackage{bbding}
\usepackage{manfnt}
\usepackage{adforn}
\usepackage{multicol}
\usepackage{amsmath,amssymb,amsfonts}
\usepackage{amsthm}
\usepackage[labelfont={bf,color=structurecolor}]{caption}
\captionsetup[table]{skip=3pt}
\captionsetup[figure]{skip=3pt}
\usepackage{booktabs}
\usepackage{multirow}
\usepackage{orcidlink}
\usepackage{etoolbox}

%------------------------------------------------------------------------------
% 2. Color Definitions (Blue Theme - Default)
%------------------------------------------------------------------------------
\definecolor{structurecolor}{RGB}{60,113,183}
\definecolor{main}{RGB}{0,166,82}
\definecolor{second}{RGB}{255,134,24}
\definecolor{third}{RGB}{0,174,247}
\definecolor{winered}{rgb}{0.5,0,0}
\definecolor{coverlinecolor}{RGB}{255,134,24}

% Color aliases for reference
\definecolor{structure3}{RGB}{60,113,183}
\definecolor{main3}{RGB}{0,166,82}
\definecolor{second3}{RGB}{255,134,24}
\definecolor{third3}{RGB}{0,174,247}

%------------------------------------------------------------------------------
% 3. Page Layout Adjustments (supplement to geometry in _extension.yml)
%------------------------------------------------------------------------------
% Note: Main geometry is set in _extension.yml
% Additional adjustments can be made here if needed

%------------------------------------------------------------------------------
% 4. Chapter/Section Format (titlesec) — match elegantbook.cls exactly
%------------------------------------------------------------------------------
% Chapter format - hang style (default)
% Note: elegantbook.cls uses \xchaptertitle which is "第 \thechapter{} 章"
% For Chinese chapter titles, set in \AtBeginDocument
\titleformat{\chapter}[hang]
  {\bfseries}
  {\filcenter\LARGE\bfseries\color{structurecolor} 第 \thechapter{} 章\;}
  {1pt}
  {\LARGE\bfseries\color{structurecolor}\filcenter}
  []

% Section format
\titleformat{\section}[hang]
  {\bfseries}
  {\Large\bfseries\color{structurecolor}\thesection\enspace}
  {1pt}
  {\color{structurecolor}\Large\bfseries\filright}

% Subsection format
\titleformat{\subsection}[hang]
  {\bfseries}
  {\large\bfseries\color{structurecolor}\thesubsection\enspace}
  {1pt}
  {\color{structurecolor}\large\bfseries\filright}

% Subsubsection format
\titleformat{\subsubsection}[hang]
  {\bfseries}
  {\large\bfseries\color{structurecolor}\thesubsubsection\enspace}
  {1pt}
  {\color{structurecolor}\large\bfseries\filright}

% Spacing
\titlespacing{\chapter}{0pt}{-20pt}{1.3\baselineskip}

%------------------------------------------------------------------------------
% 5. Headers and Footers (fancyhdr) — match elegantbook.cls
%------------------------------------------------------------------------------
\pagestyle{fancy}
\fancyhf{}

% Footer: centered page number
\fancyfoot[c]{\color{structurecolor}\small\thepage}

% Header: centered (chapter/section title)
\fancyhead[C]{\color{structurecolor}\normalfont\normalsize\leftmark}

% Header rule
\renewcommand{\headrule}{\color{structurecolor}\hrule width\textwidth}
\renewcommand{\headrulewidth}{1pt}

% Plain page style (chapter opening pages) — centered page number at bottom
\fancypagestyle{plain}{
  \renewcommand{\headrulewidth}{0pt}
  \fancyhf{}
  \fancyfoot[C]{\color{structurecolor}\small\thepage}
  \renewcommand{\headrule}{}
}

% Section and chapter marks
\renewcommand{\sectionmark}[1]{\markright{\thesection\ #1}}
\renewcommand{\chaptermark}[1]{\markboth{第 \thechapter{} 章\ #1}{}}

%------------------------------------------------------------------------------
% 6. List Environments (enumitem + tikz) — match elegantbook.cls
%------------------------------------------------------------------------------
\setlist{nolistsep}

% Itemize custom bullets
\newcommand*{\eitemi}{\tikz[baseline]\draw[baseline, ball color=structurecolor,draw=none] circle (2pt);}
\newcommand*{\eitemii}{\tikz[baseline]\draw[baseline, fill=structurecolor,draw=none,circular drop shadow] circle (2pt);}
\newcommand*{\eitemiii}{\tikz[baseline]\draw[baseline, fill=structurecolor,draw=none] circle (2pt);}

\setlist[itemize,1]{label={\eitemi}}
\setlist[itemize,2]{label={\eitemii}}
\setlist[itemize,3]{label={\eitemiii}}

% Enumerate format
\setlist[enumerate,1]{label=\color{structurecolor}\arabic*.}
\setlist[enumerate,2]{label=\color{structurecolor}(\alph*).}
\setlist[enumerate,3]{label=\color{structurecolor}\Roman*.}
\setlist[enumerate,4]{label=\color{structurecolor}\Alph*.}

%------------------------------------------------------------------------------
% 7. Theorem Environments (tcolorbox - fancy mode, match elegantbook.cls)
%------------------------------------------------------------------------------
% We define the tcolorbox styles here (same as elegantbook.cls),
% but defer the actual \newenvironment definitions to \AtBeginDocument
% (after Quarto's amsthm \newtheorem) to avoid conflicts.

\tcbuselibrary{theorems,breakable}

% Common style for all theorem boxes — match elegantbook.cls exactly
% Note: cls uses \citshape; for Chinese we use \kaishu via \citshape
% We'll define \citshape if not already defined (ctex provides \kaishu)
\providecommand{\citshape}{\itshape}

\tcbset{
  common/.style={
    fontupper=\citshape,
    lower separated=false,
    coltitle=white,
    colback=gray!5,
    boxrule=0.5pt,
    fonttitle=\bfseries,
    enhanced,
    breakable,
    top=8pt,
    before skip=8pt,
    attach boxed title to top left={yshift=-0.11in,xshift=0.15in},
    boxed title style={boxrule=0pt,colframe=white,arc=0pt,outer arc=0pt},
    separator sign={.},
  }
}

% defstyle (definition-like environments) — match cls exactly
\tcbset{
  defstyle/.style={
    common,
    colframe=main,
    colback=main!5,
    colbacktitle=main,
    overlay unbroken and last={
      \node[anchor=south east, outer sep=0pt] at (\linewidth-width,0) 
      {\textcolor{main}{$\clubsuit$}};
    },
  }
}

% thmstyle (theorem-like environments) — match cls exactly
\tcbset{
  thmstyle/.style={
    common,
    colframe=second,
    colback=second!5,
    colbacktitle=second,
    overlay unbroken and last={
      \node[anchor=south east, outer sep=0pt] at (\linewidth-width,0) 
      {\textcolor{second}{$\heartsuit$}};
    },
  }
}

% prostyle (proposition-like environments) — match cls exactly
\tcbset{
  prostyle/.style={
    common,
    colframe=third,
    colback=third!5,
    colbacktitle=third,
    overlay unbroken and last={
      \node[anchor=south east, outer sep=0pt] at (\linewidth-width,0) 
      {\textcolor{third}{$\spadesuit$}};
    },
  }
}

%------------------------------------------------------------------------------
% 7b. Theorem environment definitions moved to before-body.tex
%     (to run after Quarto's amsthm \newtheorem, avoiding conflicts)
%------------------------------------------------------------------------------

%------------------------------------------------------------------------------
% 8. Other Environments (example, exercise, problem, note, proof, solution, etc.)
%    Match elegantbook.cls exactly
%------------------------------------------------------------------------------
% Example environment
\newcounter{example}[chapter]
\renewcommand{\theexample}{\thechapter.\arabic{example}}
\newenvironment{example}[1][]{
  \refstepcounter{example}
  \par\noindent\textbf{\color{main}示例 \theexample #1 }\rmfamily}{
  \par\ignorespacesafterend}

% Exercise environment
\newcounter{exercise}[chapter]
\renewcommand{\theexercise}{\thechapter.\arabic{exercise}}
\newenvironment{exercise}[1][]{
  \refstepcounter{exercise}
  \par\noindent\makebox[-3pt][r]{\scriptsize\color{red!90}\HandPencilLeft\quad}
  \textbf{\color{main}练习 \theexercise #1 }\rmfamily}{
  \par\ignorespacesafterend}

% Problem environment
\newcounter{problem}[chapter]
\renewcommand{\theproblem}{\thechapter.\arabic{problem}}
\newenvironment{problem}[1][]{
  \refstepcounter{problem}
  \par\noindent\textbf{\color{main}问题 \theproblem #1 }\rmfamily}{
  \par\ignorespacesafterend}

% Note environment
\newenvironment{note}{
  \par\noindent\makebox[-3pt][r]{\scriptsize\color{red!90}\textdbend\quad}
  \textbf{\color{second}注} \citshape}{\par}

% proof, solution, remark environments are defined in before-body.tex
% (to run after Quarto's amsthm definitions)

% Assumption environment
\newenvironment{assumption}{\par\noindent\textbf{\color{third}假设} \citshape}{\par}

% Conclusion environment
\newenvironment{conclusion}{\par\noindent\textbf{\color{third}结论} \citshape}{\par}

% Property environment
\newenvironment{property}{\par\noindent\textbf{\color{third}性质} \citshape}{\par}

% Custom environment (user-defined name)
\newenvironment{custom}[1]{\par\noindent\textbf{\color{third} #1} \citshape}{\par}

%------------------------------------------------------------------------------
% 9. Introduction and Problemset Environments — match elegantbook.cls
%------------------------------------------------------------------------------

% Introduction environment style
\tcbset{
  introductionsty/.style={
    enhanced,
    breakable,
    colback=structurecolor!10,
    colframe=structurecolor,
    fonttitle=\bfseries,
    colbacktitle=structurecolor,
    fontupper=\citshape,
    attach boxed title to top center={yshift=-3mm,yshifttext=-1mm},
    boxrule=0pt,
    toprule=0.5pt,
    bottomrule=0.5pt,
    top=8pt,
    before skip=8pt,
    sharp corners
  }
}

% Introduction environment
% Note: Pandoc generates \begin{itemize}\tightlist\item...\end{itemize}
% inside the div. We wrap with tcolorbox + multicols but let the
% Pandoc-generated itemize provide the list structure.
\newenvironment{introduction}[1][内容提要]{
  \begin{tcolorbox}[introductionsty,title={#1}]
    \begin{multicols}{2}
      \renewcommand{\labelitemi}{\textcolor{structurecolor}{\upshape\scriptsize\SquareShadowBottomRight}}}{
    \end{multicols}
  \end{tcolorbox}}

% Problemset environment
% Pandoc generates \begin{enumerate}\tightlist\item...\end{enumerate}
% inside the div. We just provide the decorative header.
\newenvironment{problemset}[1][习题]{
  \vspace*{10pt}
  \begin{center}
    \textcolor{structurecolor}{\Large\bfseries\adftripleflourishleft~#1~\adftripleflourishright}
  \end{center}}{
  \par}

%------------------------------------------------------------------------------
% 10. Hyperref Settings (supplement to Quarto's default)
%------------------------------------------------------------------------------
% Note: colorlinks, linkcolor, citecolor, urlcolor are set in _extension.yml
% Additional hyperref settings

%------------------------------------------------------------------------------
% 11. Chinese Font Settings (for Chinese documents) — match elegantbook.cls ctexfont option
%------------------------------------------------------------------------------
% Requires XeLaTeX engine
% Use ctex with default fontset (auto-detect system fonts)
\usepackage[UTF8,scheme=plain]{ctex}
\xeCJKsetup{AutoFakeBold=true}
% Set Chinese fonts for macOS (adjust for other platforms)
%\setCJKmainfont{Songti SC}
%\setCJKsansfont{Heiti SC}
%\setCJKmonofont[AutoFakeBold=true]{FangSong_GB2312}

% Define \citshape (楷体 = \kaishu) if ctex is loaded
\ifdefined\kaishu
  \renewcommand{\citshape}{\kaishu}
\else
  \renewcommand{\citshape}{\itshape}
\fi
% Define \cfs (仿宋 = \fangsong) if ctex is loaded
\ifdefined\fangsong
  \newcommand{\cfs}{\fangsong}
\else
  \newcommand{\cfs}{\normalfont}
\fi

%------------------------------------------------------------------------------
% 12. Additional Styling Adjustments — match elegantbook.cls
%------------------------------------------------------------------------------

% Display math spacing
\setlength{\abovedisplayskip}{3pt}
\setlength{\belowdisplayskip}{3pt}

% Caption setup (already done above, repeated for clarity)
\captionsetup[table]{skip=3pt}
\captionsetup[figure]{skip=3pt}

% Listings style (code blocks) — match elegantbook.cls
\RequirePackage{listings}
\lstset{
  basicstyle=\ttfamily\small,
  breaklines=true,
  frame=single,
  rulecolor=\color{structurecolor},
  framerule=0.2pt,
  numbers=none,
  keywordstyle=\color{winered},
  commentstyle=\color{gray},
}

%------------------------------------------------------------------------------
% 13. Title Page — match elegantbook.cls \maketitle exactly
%------------------------------------------------------------------------------

% Chinese title page labels (from elegantbook.cls lang=cn)
\newcommand{\authorname}{\citshape 作者：}
\newcommand{\institutename}{\citshape 组织：}
\newcommand{\datename}{\citshape 时间：}
\newcommand{\versionname}{\citshape 版本：}

% Define ElegantBook-specific commands
\makeatletter
\newcommand{\subtitle}[1]{\gdef\@subtitle{#1}}
\newcommand{\institute}[1]{\gdef\@institute{#1}}
\newcommand{\version}[1]{\gdef\@version{#1}}
\newcommand{\extrainfo}[1]{\gdef\@extrainfo{#1}}
\newcommand{\logo}[1]{\gdef\@logo{#1}}
\newcommand{\cover}[1]{\gdef\@cover{#1}}
\newcommand\bioinfo[2]{\gdef\@bioinfo{{\citshape #1}：#2}}
\makeatother

% \cnormal is used in title page author/institute lines
\newcommand{\cnormal}{\kaishu}

% ---- Set cover and logo (defaults; override in project) ----
% These are set here so \maketitle can use them immediately.
% Users can override by redefining \@cover / \@logo before \maketitle.
\makeatletter
\gdef\@cover{images/cover.jpg}%
\gdef\@logo{images/logo-blue.png}%
\makeatother

% ---- Redefine \maketitle to match elegantbook.cls ----
% Must be in preamble (NOT \AtBeginDocument) so it overrides
% the standard \maketitle before \frontmatter\maketitle executes.
%
% The \@->\spacefactor error is caused by parskip.sty (loaded by
% Quarto's default template) which tries to process paragraph
% transitions after \includegraphics while still in vertical mode.
% Fix: wrap \includegraphics in \mbox{} to keep it in horizontal mode,
% then use explicit \par to transition back.
\makeatletter
\renewcommand*{\maketitle}{%
  \hypersetup{pageanchor=false}
  \pagenumbering{Alph}
  \begin{titlepage}
    \newgeometry{margin=0in}
    \parindent=0pt
    \parskip=0pt
    % Cover image — wrapped in \mbox to prevent
    % parskip \@ expansion in vertical mode
    \noindent\mbox{\includegraphics[width=\linewidth]{\@cover}}\par
    \nointerlineskip
    \setlength{\fboxsep}{0pt}%
    \noindent\colorbox{coverlinecolor}{\rule{0pt}{0.5in}\hspace{\linewidth}}\par
    \vfill
    \vskip-2ex
    \hspace{2em}\parbox{0.8\textwidth}{%
      \bfseries\Huge\@title\par}%
    \vfill
    \vspace{-1.0cm}
    \setstretch{2.5}
    \hspace{2.5em}%
    \begin{minipage}[c]{0.67\linewidth}
      {\color{darkgray}\bfseries\Large
        \ifcsname @subtitle\endcsname\@subtitle\\[2ex]\fi}
      \color{gray}\normalsize
      {\renewcommand{\arraystretch}{0.618}
      \begin{tabular}{l}
        \ifx\@author\empty\else\authorname\cnormal\@author\\ \fi
        \ifcsname @institute\endcsname \institutename \cnormal\@institute\\ \fi
        \datename\cnormal\zhtoday \\
        \ifcsname @version\endcsname \cnormal\versionname\@version\\ \fi
        \ifcsname @bioinfo\endcsname \cnormal\@bioinfo\\ \fi
      \end{tabular}}%
    \end{minipage}%
    \begin{minipage}[c]{0.27\linewidth}
    \begin{tikzpicture}[remember picture,overlay]
      \begin{pgfonlayer}{background}
        \node[opacity=0.8,
              anchor=south east,
              outer sep=0pt,
              inner sep=0pt] at ($(current page.south east) +(-0.8in,1.5in)$) {
                \ifcsname @logo\endcsname\includegraphics[width=4.2cm]{\@logo}\fi};
      \end{pgfonlayer}
    \end{tikzpicture}
    \end{minipage}
    \vfill
    \begin{center}
      \setstretch{1.3}
      \parbox[t]{0.7\textwidth}{\centering\citshape
        \ifcsname @extrainfo\endcsname\@extrainfo\fi}
    \end{center}
    \vfill
  \end{titlepage}
  \restoregeometry
  \thispagestyle{empty}%
  \pagenumbering{Roman}% switch to uppercase Roman numerals for TOC pages
}
\makeatother

%------------------------------------------------------------------------------
% 14. natbib citation style — numbered, sorted & compressed, square brackets
%------------------------------------------------------------------------------
\AtBeginDocument{%
  \setcitestyle{numbers,sort&compress,square}%
}

%------------------------------------------------------------------------------
% 15. Equation numbering — number within chapter (e.g. (1.1), (1.2), …)
%------------------------------------------------------------------------------
\numberwithin{equation}{chapter}

%------------------------------------------------------------------------------
% 16. Chapters start on odd pages; redefined \cleardoublepage to force
%     odd pages even with oneside (standard \cleardoublepage only does
%     \clearpage in oneside mode, which doesn't ensure odd page numbers).
%------------------------------------------------------------------------------
\makeatletter
\def\cleardoublepage{\clearpage\ifodd\c@page\else
    \hbox{}\thispagestyle{empty}\newpage\fi}
\makeatother

\AtBeginDocument{%
  \let\oldchapter\chapter
  \renewcommand{\chapter}{%
    \cleardoublepage
    \oldchapter
  }%
}%

%------------------------------------------------------------------------------
% 17. Equation reference parentheses — Quarto 对公式交叉引用使用 \ref
%     而非 \eqref，导致编号不带括号。AtBeginDocument 在 hyperref 加载后执行，
%     确保 \ref 重定义生效（覆盖 hyperref 版本）。
%------------------------------------------------------------------------------
\usepackage{xstring}
\AtBeginDocument{%
  \let\quartoOldRef\ref
  \renewcommand{\ref}[1]{%
    \IfBeginWith{#1}{eq-}{%
      \textup{(\quartoOldRef{#1})}%
    }{%
      \quartoOldRef{#1}%
    }%
  }%
}

%==============================================================================
% End of ElegantBook PDF Style Preamble (all other hooks cleaned up)
%==============================================================================
